\begin{section}{Einleitung}

Betriebssysteme bilden das Fundament moderner Computer.
Sie koordinieren Ressourcen, stellen Schnittstellen bereit und ermöglichen es, Anwendungen auf einer
einheitlichen Basis auszuführen.
Gleichzeitig sind sie traditionell ein Forschungs- und Lehrfeld, um grundlegende Konzepte der Systemsoftware zu erproben.
Insbesondere minimalistische Betriebssysteme erlauben es, Funktionen von Grund auf zu verstehen und gezielt zu erweitern.

In diesem Seminarprojekt haben wir uns mit MOROS (Minimal Obscure Rust Operating System) beschäftigt.
MOROS ist ein in \texttt{Rust} geschriebenes, bewusst
klein gehaltenes Betriebssystem, das rein textbasiert arbeitet und zunächst nur über elementare Komponenten wie Shell, Dateisystem und Netzwerkstack verfügt.
Sein reduziertes Design macht MOROS zu einer interessanten Plattform, um gezielt neue Funktionalitäten zu erproben.

Unser Ziel war es, MOROS um eine Anwendung zu erweitern, die einerseits technische Herausforderungen mit sich bringt, andererseits aber auch anschaulich und
unterhaltsam demonstriert, was mit einem minimalistischen Betriebssystem möglich ist.
Dazu haben wir ein verteiltes Mehrspieler-Spiel entwickelt, das grundlegende
Erweiterungen auf mehreren Ebenen erforderte: von der Implementierung einer Maus-Schnittstelle über eine einfache Grafik-Ausgabe bis hin zur Kommunikation über
das Netzwerk.

Die folgende Ausarbeitung beschreibt dieses Projekt im Detail.
Zunächst wird MOROS selbst vorgestellt und erläutert, welche zusätzlichen Komponenten notwendig
waren, um die Spieleentwicklung zu ermöglichen.
Anschließend wird das Spiel mit seinen Kernmechanismen wie Map-Generierung, Spieler- und Bullet-Bewegung
sowie Rendering beschrieben.
Darauf folgt die Verteilung über das Netzwerk und die zugehörige Server-Client-Architektur.
Abschließend erfolgt eine Analyse und
Bewertung des Projekts, bevor ein Fazit und mögliche Perspektiven aufgezeigt werden.

\end{section}
